% Auto-converted from Optimizing_Gestural_Efficiency.docx via pandoc
% Source section: Case Study and Representative
Application
\chapter{Case Study and Representative Application}\label{case-study-and-representative-application}

To facilitate easily-learned, gesture-efficient, on-demand non-qwerty
input, \emph{\textbf{µLex}} input applications should enable:

\begin{itemize}
\item
  Easy switching between qwerty and non-qwerty.
\item
  `Touch typing' non-qwerty with high Gestural Efficiency.
\item
  Accessibility via keyboards, mice, and other input peripherals
\item
  Qwerty access to non-qwerty characters
\end{itemize}

In the most general statement, and because the keyboard is among the
most ubiquitous input peripherals, the non-qwerty input problem reduces
to presenting \uline{qwerty-accessible} choices that result in the
highest number of non-qwerty characters for the fewest gestures. The
general strategy for optimizing the solution is to present choice sets
of multi-character non-qwerty lexemes - rather than single characters -
for selection.

To initiate discussion, consider the example of laptop users in the
field of linguistics writing papers in English about aspects of Chinese.
For this example, most user time and effort are spent on word processor
GUIs occupying monitors (Figure 15, left) and most of the document text
is English, input using laptop keyboards. Periodically, and for
relatively short periods of time, users input Chinese instead of
English. For this use case, the non-qwerty input set is Mandarin
characters and lexemes.

Three resources are used to implement the \emph{\textbf{µLex}} strategy:

\begin{itemize}
\item
  A set of Chinese characters, optionally with frequencies
\item
  A Chinese lexeme frequency list
\item
  A method for mapping qwerty characters to Chinese characters and
  lexemes
\end{itemize}

\section{\texorpdfstring{ Chinese Character (Frequency)
Lists}{ Chinese Character (Frequency) Lists}}\label{chinese-character-frequency-lists}

Given the productive nature of language, no Chinese lexicon includes
every possible lexeme. To enter lexemes not in the source frequency
list, \textbf{\emph{µ}短语}enables input of individual \textbf{汉字}
(`hànzì', \textbackslash Chinese characters\textbackslash). The
character frequency lists used in this research are from (Da, 2004). For
completeness, an additional list without frequencies that has all
individual Unicodes with a pīnyīn pronunciation is generated
programmatically.

To indicate available non-lexeme \textbf{汉字}, \textbf{\emph{µ}短语}
suffixes the green corner display text with a `+'.

For this research, the Unicode standard is used as the resource for the
complete set of Mandarin characters.

\section{Chinese Lexemes Frequency Lists From Internet Corpora
Research}\label{chinese-lexemes-frequency-lists-from-internet-corpora-research}

Lexeme frequency lists can be made using readily available tools.
(Baroni \& Kilgarriff, 2006). The two Mandarin frequency lists used for
this research are \href{http://corpus.leeds.ac.uk/frqc/lcmc.num}{The
Lancaster Corpus of Mandarin Chinese} (\textbf{LCMC},
\textasciitilde5,000 lexemes) (McEnery \& Xiao, 2004) and
\href{http://corpus.leeds.ac.uk/frqc/internet-zh.num}{Chinese Internet
Corpus} (\textbf{CIC}, \textasciitilde50,000 lexemes) (Sharoff, 2006).

\section{Pīnyīn Maps QWERTY to Chinese Character
Pronunciation}\label{pux12bnyux12bn-maps-qwerty-to-chinese-character-pronunciation}

The most widely used system for representing Chinese pronunciation is
the pīnyīn system, all of the non-tone characters for which are on the
qwerty keyboard (Her, 2005).

\section{\texorpdfstring{Example \emph{µLex} App for Chinese Lexical
Input}{Example µLex App for Chinese Lexical Input}}\label{example-uxb5lex-app-for-chinese-lexical-input}

\textbf{\emph{µ}短语} (mu-duǎnyǔ) - a family of \emph{\textbf{µLex}}
Chinese lexical input applications \textbf{--} vary with the frequency
lists used to generate them. Figure 14 shows the initial grid display
for a `small' version based on LCMC which is used as the basis for this
research. A larger version, based on CIC, is used for comparison
purposes.

\begin{figure}[htbp]
\centering
\includegraphics[width=4.81042in,height=1.37361in]{media/media/image10.emf}
\caption{\emph{µ}\textbf{短语} Initial Supplemental Display for Chinese Lexical Input}
\label{fig:auto:1}
\end{figure}


\section{\texorpdfstring{Enabling and Disabling
\textbf{\emph{µ}短语} with \emph{Long-press}
}{Enabling and Disabling µ短语 with Long-press }}\label{enabling-and-disabling-uxb5ux77edux8bed-with-long-press}

For the example use case, an always-visible \textbf{\emph{µ}短语} grid
is unnecessary and potentially distracting. So, \textbf{\emph{µ}短语}
increases usability by making it easy to \textbf{\uline{enable}} and
\textbf{\uline{disable}} (\emph{~}Figure 14 and Figure 15, red ellipse).
Enabling \textbf{\emph{µ}短语} is done via \emph{\textbf{activation}} or
\emph{\textbf{toggling}}. If it is being used for the first time since
boot, users \textbf{\uline{activate}} by selecting an item in an OS
menu. If the software is already \textbf{\uline{active}},
\emph{\textbf{long-press}} of the \textbf{\textless/\textgreater{}} key
\textbf{\uline{toggles it on}} (\emph{\textbf{show}}) and
\textbf{\uline{toggles it off}} (\emph{\textbf{hide}}).
\textbf{\uline{Deactivation}} is via mouse-clickable selections accessed
through the blue corner cell. Since \textbf{activation} gesture is via
the OS, \textbf{show} and \textbf{hide} are via
\emph{\textbf{long-press}}, and \textbf{deactivate} is via
\textbf{\emph{µ}短语}'s grid, none of these interfere with the GFAs of
concurrently running applications.

\begin{figure}[htbp]
\centering
\includegraphics[width=6.24638in,height=1.7963in]{media/media/image11.emf}
\caption{\textbf{\emph{µ}短语} Supplements Other Apps to Provide Chinese Lexical Input}
\label{fig:auto:2}
\end{figure}


\section{\texorpdfstring{Using \textbf{\emph{µ}短语} for Selecting
Lexemes}{Using µ短语 for Selecting Lexemes}}\label{using-uxb5ux77edux8bed-for-selecting-lexemes}

\paragraph{Select the Most Frequently Used Lexemes with 1
Gesture}\label{select-the-most-frequently-used-lexemes-with-1-gesture}

The initial \emph{µ}\textbf{短语} display shows the most frequently used
lexemes according to the source corpus. When the grid is visible, users
select lexemes by mouse-clicking their cells or by
\emph{\textbf{long-pressing}} keys with marks indicated in blue (Figure
14, red circle). For example, typing \textbf{\textless K\textgreater{}}
in the usual manner inputs lowercase a and \emph{\textbf{long-pressing}}
\textbf{\textless K\textgreater{}} inputs lexeme\textbf{这}
(\textbackslash zhè\textbackslash). By using long-press,
\emph{µ}\textbf{短语} does not interfere with existing word-processing
GFAs. Feedback via highlight (Figure 14, yellow) indicates crossing the
\emph{\textbf{long-press}} interval threshold.

Another example: to input \textbf{的} (`de`,
\textbackslash possessive\textbackslash), users press
\textbf{\textless A\textgreater{}} until the cell highlights, and it is
input upon release of the key.

The GE for any single-character lexeme in the \textbf{\emph{µ}短语}grid
is 1.00 and the GE for \textbf{一个} (`yīgè') is 2.00. This contrasts
with the Microsoft Pinyin IME where the number of gestures to input
\textbf{的}varies from 2 (\textbf{d},
\textbf{\textless enter\textgreater{}}) to 3 (\textbf{d}, \textbf{e},
\textbf{1}) with corresponding GEs of 0.50 and 0.33, not including the
gestures necessary to switch between English and Chinese. The lexemes on
the initial display account for 18.69\% of written Chinese according to
LCMC.

\paragraph{\texorpdfstring{Select the Next Most Frequently Used Lexemes
with 2 Gestures
}{Select the Next Most Frequently Used Lexemes with 2 Gestures }}\label{select-the-next-most-frequently-used-lexemes-with-2-gestures}

To input lexemes not on the initial display, users
\textbf{\textless shift\textgreater{}} \emph{\textbf{long-press}} the
lexeme's first pīnyīn. So, to input \textbf{时候 (}`shíhòu'
\textbackslash when\textbackslash), users:

\begin{itemize}
\item
  \textbf{\textless shift\textgreater{} \emph{Long-press}}
  \textbf{\textless S\textgreater{}} to update the display (Figure 3)
\item
  \emph{\textbf{Long-press}} \textbf{\textless P\textgreater{}} to input
  the lexeme
\end{itemize}

\begin{figure}[htbp]
\centering
\includegraphics[width=4.83333in,height=1.41875in]{media/media/image12.emf}
\caption{: Grid of Most Frequently Used Lexemes with Pronunciations Beginning with S}
\label{fig:auto:3}
\end{figure}


There are 658 frequently used lexemes that can be similarly selected
with 2 gestures. The average GE for all 2-gesture lexemes is 0.80. and
2-gesture lexemes account for 25.78\% of written Chinese. So, 44.47\% of
written Chinese can be input with 2 or fewer gestures.

\paragraph{66\% of Written Chinese in 4 or Fewer
Gestures}\label{of-written-chinese-in-4-or-fewer-gestures}

\includegraphics[width=4.21in,height=2.29in]{media/media/image13.emf}

\begin{enumerate}
\def\labelenumi{\arabic{enumi}.}
\setcounter{enumi}{3}
\tightlist
\item
\end{enumerate}
